\scanpage{42}
Soit $S$ un schéma de base, et soit $X$ un schéma de Br. Severi sur $S$ muni d'une \uline{\unsure{trisection}} non ramifiée $X'$, i.e. un ss-schéma \inserted{fermé} de $X$ qui est non ramifié sur $X$, d'ordre 3 sur $X$. Un tel truc peut être considéré comme un objet à \unsure{schém.} structural $\mathfrak S_3$, opérant sur $P^1_S$ comme il a été dit. La classification de tels objets est donc la même que la classification dans un $H^1(S,\mathfrak S_3)$ des revêtts \inserted{gal.} \uline{non ramifiés} de groupe $\mathfrak S_3$ de $S$. De plus, $\sigma$ définit dans $X$ un $S$-automorphisme \inserted{involutif} \unsure{intr.} $\unsure{\sigma_S|\sigma}$ qui transforme $X'$ en un $X''$, dit \uline{associé} à $X'$. Évidt, si on partait de $(X,X'')$, on aurait le même $\sigma$, et \unsure{on restitue} $X'$ [à préciser\ldots]. De plus, on obtient [par descente\ldots] une \uline{bisection}, non ramifiée \deleted{\unsure{[en ses points}} en les pts de $S$ de car. $\ne2,3$]..... On peut passer au quotient, \deleted{\unsure{f}} par $\mathfrak S_3$, \deleted{\unsure{puis}} on trouve \unsure{par} :
\[
P^1_S\xrightarrow{p}P^1_S\quad\text{(comp. car }\mathfrak S_3\text{ opère}
\]
\mspage{41}
